\clearpage 

\section {Summary }
\label{sec:summary}

In this note, we have presented the results of a 2-flavour $\nu_{\mu}$-disappearance analysis 
on the combined Run 1+2 dataset, which corresponds to an integrated JPARC neutrino beam exposure 
of 1.431 $\times$ 10$^{20}$ POT. The analysis method was identical to that used in our earlier analysis
reported in Ref.\ \cite{OurRun1DisappAnalysis} except that updates were used of the SuperK flux, 
the flux normalisation error ($N_{SK}/N_{ND}$), the flux shape error, the ND280 event rate 
normalisation correction, the SuperK ring-counting efficiency error and the SuperK energy scale error.

The ND280 reconstructed energy spectrum was extrapolated to predict the number of single $\mu$-like 
ring events that would be expected in SuperK if there were no oscillations. With GENIE this prediction was 
101.72 $\pm$ 10.09 (stat) $^{+13.88}_{-13.49}$ (syst), while NEUT predicted 103.59 $\pm$ 10.18 (stat) 
$^{+13.77}_{-13.43}$ (syst) events. Only 31 single $\mu$-like ring events were observed in SuperK, 
and the statistical significance of the deficit was $\sim$4.2 standard deviations for GENIE and 
$\sim$4.3 standard deviations for NEUT. 

A 2-flavour $\nu_{\mu}$-disappearance fit to the observed SuperK reconstructed energy spectrum of 
single $\mu$-like ring events in the combined Run 1+2 dataset gave best-fit values of 
sin$^{2}$(2$\theta_{23}$) = 0.97 and $|\Delta m^{2}_{23}|$ = 2.60 $\times$ 10$^{-3}$ eV$^{2}$/c$^{4}$
($\chi^{2}$/ndf = 0.310) for GENIE and sin$^{2}$(2$\theta_{23}$) = 0.98 and $|\Delta m^{2}_{23}|$ 
= 2.65 $\times$ 10$^{-3}$ eV$^{2}$/c$^{4}$ ($\chi^{2}$/ndf = 0.311) for NEUT. Separate 2-flavour 
$\nu_{\mu}$-disappearance fits were also made to each of the Run 1 and Run 2 datasets. For Run 1
these gave best-fit values of sin$^{2}$(2$\theta_{23}$) = 0.81 and $|\Delta m^{2}_{23}|$ 
= 2.77 $\times$ 10$^{-3}$ eV$^{2}$/c$^{4}$ ($\chi^{2}$/ndf = 0.105) for GENIE and 
sin$^{2}$(2$\theta_{23}$) = 0.84 and $|\Delta m^{2}_{23}|$ = 2.88 $\times$ 10$^{-3}$ eV$^{2}$/c$^{4}$ 
($\chi^{2}$/ndf = 0.114) for NEUT. Meanwhile the best-fit values for Run 2 were 
sin$^{2}$(2$\theta_{23}$) = 1.00 and $|\Delta m^{2}_{23}|$ = 2.59 $\times$ 10$^{-3}$ eV$^{2}$/c$^{4}$ 
($\chi^{2}$/ndf = 0.318) for GENIE and sin$^{2}$(2$\theta_{23}$) = 1.00 and $|\Delta m^{2}_{23}|$ 
= 2.62 $\times$ 10$^{-3}$ eV$^{2}$/c$^{4}$ ($\chi^{2}$/ndf = 0.310) for NEUT.

Confidence contours for the oscillation parameters sin$^{2}$(2$\theta_{23}$) and 
$|\Delta m^{2}_{23}|$ were constructed using the Feldman-Cousins method. These contours have been 
presented at 68\% and 90\% confidence both with and without systematic errors for the combined 
Run 1+2 dataset. Contours have also been shown at 90\% confidence with systematic errors 
for each of Run 1 and Run 2 separately. The 68\% and 90\% confidence contours with and without 
systematic errors for the combined Run 1+2 dataset have been compared with previous results 
from MINOS and SuperK. 




